\begin{tcolorbox}[colback=red!3!white,colframe=red!50!black,boxrule=0.8pt,title={Correction notes},fonttitle=\bfseries]
\textbf{Corrections from the calculation check (30~Sept.~2026).} The following places are corrected or annotated in the text; each carries a dated note with the previous value:
\begin{itemize}
\item EM geometry factor $\sqrt{4\pi/9} = 1.18270 \to 1.18164$; $\xi_\ell = 2E_P/(\lambda_C E_P) \to 2\ell_P/\bar{\lambda}_C$; note on $\xi_E = 2\sqrt{G}\,E$ with values $10^5$--$10^9$ (at most $2$ in natural units).
\item (1~Oct.~2026) Higgs expression $\lambda_h^2 v^2/(16\pi^3 E_h^2)$: $1.316$ and $1.32\times10^{-4} \to 1.30\times10^{-4}$ (with $\lambda_h = m_h^2/(2v^2) \approx 0.129$); the comparison value ``CODATA/Experimental'' $1.320\times10^{-4}$ for $\xi$ does not exist and is replaced by the setting $4/30000$; agreement $99.7\,\% \to 97.6\,\%$; ``$= \xi$'' $\to$ ``$\approx \xi$'' (consistency check, Doc.~385).
\item (2~Oct.~2026) Notation $f$: general dependence $m_e = f(\xi, \text{Higgs})$ --- not a frequency; note in the text (Doc.~190, R147).
\end{itemize}
\end{tcolorbox}

\section{Introduction: Ratio-Based vs. Parameter-Based Physics}

	\medskip\noindent\textbf{Note on versions.} The German and English versions of this document differ in structure and extent; earlier revisions shortened, extended or changed parts of both. Where they differ, the more extensive German version applies; sections contained only in this English version count as supplements.

	
	This document presents a verification of the T0 model, starting from the view that $\xi$ is a scale ratio, not an assigned numerical value. This distinction is important for understanding the parametrization of the T0 model, which works with the single parameter $\xi$ (plus integer structural choices as motivated settings; SI anchors serve for conversion).
	
	\medskip\noindent\textbf{Status markers.} Statements marked \textbf{[S]} are \emph{speculative}: a hint or conjecture that has not yet been derived or numerically checked in the corpus.
	
	\begin{tcolorbox}[colback=red!5!white,colframe=red!75!black,title=Two Formulations]
		\textbf{Usual practice (literature):}
		\begin{align}
			\xi &= 1.32 \times 10^{-4} \quad \text{(numerical value assigned)} \\
			\alpha_{\text{EM}} &= \frac{1}{137} \quad \text{(numerical value assigned)} \\
			G &= 6.67 \times 10^{-11} \quad \text{(numerical value assigned)}
		\end{align}
		
		\textbf{T0 formulation:}
		\begin{align}
			\xi &\approx \frac{\lambda_h^2 v^2}{16\pi^3 E_h^2} \quad \text{(Higgs energy scale ratio, consistency check)} \\
			\xi &= \frac{2\ell_P}{\lambda_C} \quad \text{(Planck-Compton length ratio)}
		\end{align}
		\textit{(Corrected on 1~Oct.~2026: previously $\xi = \lambda_h^2 v^2/(16\pi^3 E_h^2)$ -- with $\lambda_h = m_h^2/(2v^2) \approx 0.129$ the Higgs expression gives $\xi \approx 1.30\times10^{-4}$, about $2.3\,\%$ below $4/30000$; a consistency check, not an exact derivation, Doc.~385.)}
	\end{tcolorbox}
	
	\section{Calculation Verification}
	
	The following tables compare T0 calculations based on scale ratios with established SI reference values. All tables are scaled to fit Kindle/portrait format. Column \emph{Status}: yes = agrees, new = new prediction without measured value, no = does not agree.
	
	% TABLE 1: Fundamental Constants and Derived Quantities
	\begin{table}[H]
		\centering
		\caption{T0 Model Verification – Part 1: Fundamental \& Derived Constants}
		\label{tab:verification_fundamental}
		\scriptsize\setlength{\tabcolsep}{3pt}\adjustbox{max width=\textwidth}{%
			\begin{tabular}{p{3.2cm}p{1.2cm}p{2.5cm}p{1.8cm}p{2.0cm}p{1.0cm}p{0.6cm}}
				\toprule
				\textbf{Physical Quantity} & \textbf{SI Unit} & \textbf{T0 Ratio Formula} & \textbf{T0 Calculation} & \textbf{CODATA/Experimental} & \textbf{Agreement} & \textbf{Status} \\
				\midrule
				$\xi$ (Flat) & 1 & $\xi \approx \frac{\lambda_h^2 v^2}{16\pi^3 E_h^2}$ & $\mathbf{1.30\times10^{-4}}$ & $4/30000$ (setting, no measured value) & $\mathbf{97.6\%}$ & yes \\
				$\xi$ (Spherical) & 1 & $\xi = \frac{\lambda_h^2 v^2}{24\pi^{5/2} E_h^2}$ & $\mathbf{1.557\times10^{-4}}$ & T0 derivation & $\mathbf{N/A}$ & new \\
				Electron mass & MeV & $m_e = f(\xi, \text{Higgs})$ & $\mathbf{0.511}$ & $0.51099895$ & $\mathbf{99.998\%}$ & yes \\
				Compton wavelength & m & $\lambda_C = \frac{\hbar}{m_e c}$ & $\mathbf{3.862\times10^{-13}}$ & $3.8615927\times10^{-13}$ & $\mathbf{99.989\%}$ & yes \\
				Planck length & m & $\ell_P$ from $\xi$ scaling & $\mathbf{1.616\times10^{-35}}$ & $1.616255\times10^{-35}$ & $\mathbf{99.984\%}$ & yes \\
				\bottomrule
			\end{tabular}
		}
	\end{table}

\textit{(Note of 2~Oct.~2026 on notation: here $m_e = f(\xi, \text{Higgs})$ is a dependence that is not written out, not a frequency and not the base winding number $f = 1/\xi$. Cf.\ Doc.~190, R147.)}

	\textit{(Corrected on 1~Oct.~2026: row $\xi$ (Flat) previously T0 calculation $1.316\times10^{-4}$, comparison value ``CODATA/Experimental'' $1.320\times10^{-4}$, agreement $99.7\,\%$ -- there is no CODATA or measured value for $\xi$, the comparison value is the setting $4/30000$; with $\lambda_h = m_h^2/(2v^2) = 0.1291$ ($m_h = 125.1$~GeV, $v = 246.22$~GeV) one has $\xi = m_h^2/(64\pi^3 v^2) = 1.301\times10^{-4}$, i.e.\ $97.6\,\%$; consistency check, Doc.~385.)}
	
	% TABLE 2: QED Corrections
	\begin{table}[H]
		\centering
		\caption{T0 Model Verification – Part 2: QED Corrections}
		\label{tab:verification_qed}
		\scriptsize\setlength{\tabcolsep}{3pt}\adjustbox{max width=\textwidth}{%
			\begin{tabular}{p{3.2cm}p{1.2cm}p{2.5cm}p{1.8cm}p{2.0cm}p{1.0cm}p{0.6cm}}
				\toprule
				\textbf{Physical Quantity} & \textbf{SI Unit} & \textbf{T0 Ratio Formula} & \textbf{T0 Calculation} & \textbf{CODATA/Experimental} & \textbf{Agreement} & \textbf{Status} \\
				\midrule
				Vertex correction & 1 & $\frac{\Delta\Gamma}{\Gamma^{\mu}} = \xi^2$ & $\mathbf{1.742\times10^{-8}}$ & New & $\mathbf{N/A}$ & new \\
				Energy indep. (1 MeV) & 1 & $f(E/E_P)$ & $\mathbf{1.000}$ & New & $\mathbf{N/A}$ & new \\
				Energy indep. (100 GeV) & 1 & $f(E/E_P)$ & $\mathbf{1.000}$ & New & $\mathbf{N/A}$ & new \\
				\bottomrule
			\end{tabular}
		}
	\end{table}
	
	% TABLE 3: Cosmological Predictions
	\begin{table}[H]
		\centering
		\caption{T0 Model Verification – Part 3: Cosmological Predictions [S]}
		\label{tab:verification_cosmological}
		\scriptsize\setlength{\tabcolsep}{3pt}\adjustbox{max width=\textwidth}{%
			\begin{tabular}{p{3.2cm}p{1.6cm}p{2.5cm}p{2cm}p{2.0cm}p{1.0cm}p{0.6cm}}
				\toprule
				\textbf{Physical Quantity} & \textbf{SI Unit} & \textbf{T0 Ratio Formula} & \textbf{T0 Calculation} & \textbf{CODATA/Experimental} & \textbf{Agreement} & \textbf{Status} \\
				\midrule
				$H_0$ (T0) & km/s/Mpc & see Doc.~026 & $\mathbf{66.2}$ & $67.4$ (Planck) & $\mathbf{1.9\%}$ & yes \\

				Universe age & Gyr & $t_U = 1/H_0$ & $\mathbf{14.0}$ & $13.8$ & $\mathbf{98.6\%}$ & yes \\
				\multicolumn{7}{l}{\textit{For complete $H_0$ derivation and measurement comparisons see Document~026.}} \\
				\bottomrule
			\end{tabular}
		}
	\end{table}
	
	\textbf{[S]} The cosmological values in Table~\ref{tab:verification_cosmological} are speculative: the present corpus deliberately leaves the cosmic sector aside, because the Standard Model does not derive it either (P39).
	
	% TABLE 4: Physical Fields and Planck Current
	\begin{table}[H]
		\centering
		\caption{T0 Model Verification – Part 4: Physical Fields \& Planck Current}
		\label{tab:verification_fields}
		\scriptsize\setlength{\tabcolsep}{3pt}\adjustbox{max width=\textwidth}{%
			\begin{tabular}{p{3.2cm}p{1.6cm}p{2.5cm}p{2cm}p{2.0cm}p{1.0cm}p{0.6cm}}
				\toprule
				\textbf{Physical Quantity} & \textbf{SI Unit} & \textbf{T0 Ratio Formula} & \textbf{T0 Calculation} & \textbf{CODATA/Experimental} & \textbf{Agreement} & \textbf{Status} \\
				\midrule
				Schwinger E-field & V/m & $E_S = \frac{m_e^2 c^3}{e\hbar}$ & $\mathbf{1.32\times10^{18}}$ & $1.32\times10^{18}$ & $\mathbf{100.0\%}$ & yes \\
				Critical B-field & T & $B_c = \frac{m_e^2 c^2}{e\hbar}$ & $\mathbf{4.41\times10^{9}}$ & $4.41\times10^{9}$ & $\mathbf{100.0\%}$ & yes \\
				Planck E-field & V/m & $E_P = \sqrt{\frac{c^7}{4\pi\varepsilon_0\hbar G^2}}$ & $\mathbf{6.45\times10^{61}}$ & $6.45\times10^{61}$ & $\mathbf{100.0\%}$ & yes \\
				Planck B-field & T & $B_P = E_P/c$ & $\mathbf{2.15\times10^{53}}$ & $2.15\times10^{53}$ & $\mathbf{100.0\%}$ & yes \\
				Planck current (Std) & A & $I_P = \sqrt{\frac{c^6\varepsilon_0}{G}}$ & $\mathbf{9.81\times10^{24}}$ & $3.479\times10^{25}$ & $\mathbf{28.2\%}$ & no \\
				Planck current (Corr) & A & $I_P = \sqrt{\frac{4\pi c^6\varepsilon_0}{G}}$ & $\mathbf{3.479\times10^{25}}$ & $3.479\times10^{25}$ & $\mathbf{99.98\%}$ & yes \\
				\bottomrule
			\end{tabular}
		}
	\end{table}
	
	\textit{(Corrected on 29~Sept.~2026: the Planck fields previously read $E_P = c^4/(4\pi\varepsilon_0 G) = 1.04\times10^{61}$~V/m and $B_P = c^3/(4\pi\varepsilon_0 G) = 3.48\times10^{52}$~T -- both formulas are dimensionally inconsistent; the correct Planck field strength is $E_P = \sqrt{c^7/(4\pi\varepsilon_0\hbar G^2)}$ with $B_P = E_P/c$.)}
	
	\section{SI-Planck Units System Verification}
	
	\subsection{Complex Formula Method vs. Simple Energy Relationships}
	
	\begin{tcolorbox}[colback=yellow!5!white,colframe=yellow!75!black,title=Observation]
		Simple relationships give numerically more accurate values than composite formulas, because fewer rounding errors accumulate.
	\end{tcolorbox}
	
	% TABLE 5: Complex Formula Method
	\begin{table}[H]
		\centering
		\caption{SI-Planck Units: Complex Formula Method}
		\label{tab:verification_complex}
		\scriptsize\setlength{\tabcolsep}{3pt}\adjustbox{max width=\textwidth}{%
			\begin{tabular}{p{3.2cm}p{1.6cm}p{2.5cm}p{2cm}p{2.0cm}p{1.0cm}p{0.6cm}}
				\toprule
				\textbf{Physical Quantity} & \textbf{SI Unit} & \textbf{Planck Formula} & \textbf{T0 Calculation} & \textbf{CODATA Reference} & \textbf{Agreement} & \textbf{Status} \\
				\midrule
				Planck time & s & $t_P = \sqrt{\frac{\hbar G}{c^5}}$ & $\mathbf{5.392\times10^{-44}}$ & $5.391\times10^{-44}$ & $\mathbf{100.016\%}$ & yes \\
				Planck length & m & $\ell_P = \sqrt{\frac{\hbar G}{c^3}}$ & $\mathbf{1.617\times10^{-35}}$ & $1.616\times10^{-35}$ & $\mathbf{100.030\%}$ & yes \\
				Planck mass & kg & $m_P = \sqrt{\frac{\hbar c}{G}}$ & $\mathbf{2.177\times10^{-8}}$ & $2.176\times10^{-8}$ & $\mathbf{100.044\%}$ & yes \\
				Planck temperature & K & $T_P = \sqrt{\frac{\hbar c^5}{G k_B^2}}$ & $\mathbf{1.417\times10^{32}}$ & $1.417\times10^{32}$ & $\mathbf{99.988\%}$ & yes \\
				Planck current & A & $I_P = \sqrt{\frac{4\pi c^6 \varepsilon_0}{G}}$ & $\mathbf{3.479\times10^{25}}$ & $3.479\times10^{25}$ & $\mathbf{99.980\%}$ & yes \\
				\bottomrule
			\end{tabular}
		}
	\end{table}
	
	\begin{tcolorbox}[colback=gray!5!white,colframe=gray!75!black,title=Note on Rounding Errors]
		Complex formulas show 99.98–100.04\% agreement due to rounding error accumulation. This is not a prediction error but a computational artifact.
	\end{tcolorbox}
	
	\subsection{Simple Energy Relationships Method}
	
	% TABLE 6: Simple Energy Relationships
	\begin{table}[H]
		\centering
		\caption{Natural Units: Simple Energy Relationships Method}
		\label{tab:verification_simple}
		\scriptsize\setlength{\tabcolsep}{3pt}\adjustbox{max width=\textwidth}{%
			\begin{tabular}{p{3.0cm}p{1.5cm}p{2.5cm}p{2.8cm}p{2.2cm}p{1.2cm}p{0.6cm}}
				\toprule
				\textbf{Physical Quantity} & \textbf{Relationship} & \textbf{Example} & \textbf{Electron Case} & \textbf{Numerical Value} & \textbf{Agreement} & \textbf{Status} \\
				\midrule
				\multicolumn{7}{l}{\textbf{DIRECT ENERGY IDENTITIES - NO ROUNDING ERRORS}} \\
				\midrule
				Mass & $E = m$ & Energy = Mass & $0.511$ MeV & Same value & $\mathbf{100\%}$ & yes \\
				Temperature & $E = T$ & Energy = Temperature & $5.93\times10^9$ K & Direct conversion & $\mathbf{100\%}$ & yes \\
				Frequency & $E = \omega$ & Energy = Frequency & $7.76\times10^{20}$ Hz & Direct identity & $\mathbf{100\%}$ & yes \\
				\midrule
				\multicolumn{7}{l}{\textbf{INVERSE ENERGY RELATIONSHIPS - EXACT}} \\
				\midrule
				Length & $E = 1/L$ & Energy = 1/Length & $3.862\times10^{-13}$ m & Inverse relationship & $\mathbf{100\%}$ & yes \\
				Time & $E = 1/T$ & Energy = 1/Time & $1.288\times10^{-21}$ s & Inverse relationship & $\mathbf{100\%}$ & yes \\
				\midrule
				\multicolumn{7}{l}{\textbf{T0 ENERGY PARAMETERS - PURE RATIOS}} \\
				\midrule
				$\xi$ (Higgs, Flat) & $E_h/E_P$ & Energy ratio & $1.30\times10^{-4}$ & Higgs check against $4/30000$ & $\mathbf{97.6\%}$ & yes \\
				$\xi$ (Higgs, Sph) & $E_h/E_P$ & Corrected ratio & $1.557\times10^{-4}$ & T0 derivation & $\mathbf{100\%}$ & new \\
				$\xi$ Geometric & $E_\ell/E_P$ & Length-energy ratio & $8.37\times10^{-23}$ & Pure geometry & $\mathbf{100\%}$ & yes \\
				EM geometry factor & Ratio & $\sqrt{4\pi/9}$ & $1.18164$ & Mathematically exact & $\mathbf{100\%}$ & new \\
				\midrule
				\multicolumn{7}{l}{\textbf{COMPLETE SI UNITS ENERGY COVERAGE - ALL 7/7 UNITS}} \\
				\midrule
				Electric current & $I = E/T$ & Energy flow rate & $[E]$ dim. & Direct energy relationship & $\mathbf{100\%}$ & yes \\
				Amount of substance & $[E^2]$ dim. & Energy density ratio & Dimensional structure & SI-defined $N_A$ & $\mathbf{Def.}$ & new \\
				Luminous intensity & $[E^3]$ dim. & Energy flow perception & Dimensional structure & SI-defined 683 lm/W & $\mathbf{Def.}$ & new \\
				\bottomrule
			\end{tabular}
		}
	\end{table}
	
	\textit{(Corrected on 30~Sept.~2026: EM geometry factor previously $1.18270$ -- $\sqrt{4\pi/9} = \sqrt{1.396263} = 1.18164$.)}
	\textit{(Corrected on 1~Oct.~2026: row $\xi$ (Higgs, Flat) previously $1.316\times10^{-4}$, ``From Higgs physics'', $100\,\%$ -- with $\lambda_h = m_h^2/(2v^2)$ one obtains $1.30\times10^{-4}$, $97.6\,\%$ of $4/30000$; consistency check.)}
	
	\begin{tcolorbox}[colback=blue!5!white,colframe=blue!75!black,title=Accuracy through Simplification]
		\textbf{Complex Formula Method (Traditional Physics):}
		\begin{itemize}
			\item Uses: $\sqrt{\frac{\hbar G}{c^5}}$, multiple constants, conversion factors
			\item Result: 99.98–100.04\% agreement (rounding errors accumulate)
			\item Problem: Each calculation step introduces small errors
		\end{itemize}
		
		\textbf{Simple Energy Relationships Method (T0 Physics):}
		\begin{itemize}
			\item Uses: Direct identities $E = m$, $E = 1/L$, $E = 1/T$
			\item Result: exact by definition (identities in natural units)
			\item Advantage: No intermediate calculations, no error accumulation
		\end{itemize}
		
		\textbf{CLASSIFICATION:}
		The difference is computational: identities in natural units avoid intermediate steps and thus rounding errors. This fits the T0 view that energy is the fundamental quantity; it is not a proof of it.
		
		\textbf{CONCLUSION}: Simple relationships are numerically no less accurate here.
	\end{tcolorbox}
	
	\section{The $\xi$ Parameter Hierarchy}
	
	\subsection{Clarification}
	
	\begin{tcolorbox}[colback=red!10!white,colframe=red!75!black,title=Note: Use of the Symbol $\xi$ in This Document]
		\textbf{Possible misunderstanding:} Treating $\xi$ as a universal parameter
		
		\textbf{In this document:} $\xi$ is a \textbf{class of dimensionless scale ratios}, not a single value.
		
		\textbf{Consequence of a mix-up:} Misinterpretations, wrong numerical values, dimensional errors.
		
		$\xi$ represents any dimensionless ratio of the form:
		\begin{equation}
			\xi = \frac{\text{T0-characteristic energy scale}}{\text{Reference energy scale}}
		\end{equation}
		
		The T0 model uses $\xi$ to denote various dimensionless ratios in different physical contexts.
	\end{tcolorbox}
	
	\subsection{The three fundamental $\xi$ energy scales}
	
	\begin{table}[H]
		\centering
		\caption{The three fundamental $\xi$ parameter types in the T0 model}
		\label{tab:xi_hierarchy}
		\scriptsize\setlength{\tabcolsep}{3pt}\adjustbox{max width=\textwidth}{%
			\begin{tabular}{|p{2.5cm}|p{4.5cm}|p{2.2cm}|p{3.8cm}|}
				\hline
				\textbf{Context} & \textbf{Definition} & \textbf{Typical Value} & \textbf{Physical Meaning} \\
				\hline
				\textbf{Energy-dependent} & $\xi_E = 2\sqrt{G} \cdot E$ & $10^5$ to $10^9$ & Energy-field coupling \\
				\hline
				\textbf{Higgs sector} & $\xi_H = \frac{\lambda_h^2 v^2}{16\pi^3 E_h^2}$ & $1.30\times10^{-4}$ & Energy scale ratio \\
				\hline
				\textbf{Scale hierarchy} & $\xi_\ell = \frac{2\ell_P}{\bar{\lambda}_C}$ & $8.37\times10^{-23}$ & Energy hierarchy ratio \\
				\hline
			\end{tabular}
		}
	\end{table}
	
	\textit{(Corrected on 1~Oct.~2026: Higgs sector previously $1.32\times10^{-4}$ -- with $\lambda_h = m_h^2/(2v^2) \approx 0.129$ one has $\xi_H = m_h^2/(64\pi^3 v^2) \approx 1.30\times10^{-4}$, about $2.3\,\%$ below $4/30000$.)}
	
	\textit{(Corrected on 30~Sept.~2026: $\xi_\ell$ previously $2E_P/(\lambda_C E_P)$ -- this reduces to $2/\lambda_C$ and is not dimensionless; the numerical value $8.37\times10^{-23}$ is $2\ell_P/\bar{\lambda}_C = 2\cdot1.616\times10^{-35}~\text{m}/3.862\times10^{-13}~\text{m}$ with $\bar{\lambda}_C = \hbar/(m_e c)$.)}
	
	\textit{(Note of 30~Sept.~2026: $\xi_E = 2\sqrt{G}\cdot E$ with a \enquote{typical value $10^5$ to $10^9$} does not hold -- in natural units $\sqrt{G} = 1/E_P$, hence $\xi_E = 2E/E_P \le 2$ for all $E \le E_P$; values $10^5$--$10^9$ are therefore excluded, and in SI $2\sqrt{G}\,E$ is not dimensionless. This also affects Rule~1 below.)}
	
	\subsection{Application rules}
	
	\begin{tcolorbox}[colback=blue!5!white,colframe=blue!75!black,title=Application Rules for $\xi$ Parameters (Pure Energy)]
		\textbf{Rule 1: Universal energy-dependent systems (RECOMMENDED)}
		\begin{equation}
			\text{Use } \xi_E = 2\sqrt{G} \cdot E \text{ where } E \text{ is the relevant energy}
		\end{equation}
		
		\textbf{Rule 2: Cosmological/coupling unification (SPECIAL CASES)}
		\begin{equation}
			\text{Use } \xi_H = 1.32 \times 10^{-4} \text{ (Higgs energy ratio)}
		\end{equation}
		
		\textbf{Rule 3: Pure energy hierarchy analysis (THEORETICAL)}
		\begin{equation}
			\text{Use } \xi_\ell = 8.37 \times 10^{-23} \text{ (energy scale ratio)}
		\end{equation}
		
		\textbf{Note:} In practice, Rule 1 applies to the vast majority of T0 calculations due to the pronounced T0 scale hierarchy.
	\end{tcolorbox}
	
	\section{Important Insights from Verification}
	
	\subsection{Main Results}
	
	\begin{tcolorbox}[colback=green!5!white,colframe=green!75!black,title=Main Results of T0 Verification]
		\textbf{1. Scale ratio validation:}
		\begin{itemize}
			\item Established values: 99.99\% agreement with CODATA
			\item Geometric $\xi$ ratio: 100.003\% agreement with Planck-Compton calculation
			\item Dimensional consistency of the checked quantities
		\end{itemize}
		
		\textbf{2. New testable predictions:}
		\begin{itemize}
			\item QED vertex ratios: $1.74 \times 10^{-8}$ (energy-independent)
			\item \textbf{[S]} Cosmological $H_0^{\text{T0}} \approx 66.2\,$km/s/Mpc (calibrated to $H_0$, see Document~026)
			\item \textbf{[S]} Redshift ratios: 40.5\% spectral variation
		\end{itemize}
		
		\textbf{3. Overall assessment:}
		\begin{itemize}
			\item Established values: 99.99\% agreement
			\item New predictions: 14+ testable ratios
			\item Dimensional consistency: given for all table quantities
			\item Scale ratio basis: internally consistent
		\end{itemize}
	\end{tcolorbox}
	
	\subsection{Experimental Testability}
	
	The ratio-based nature of the T0 model enables specific experimental tests:
	
	\begin{enumerate}
		\item \textbf{Energy scale-independent QED corrections}:
		\begin{equation}
			\frac{\Delta\Gamma^{\mu}(E_1)}{\Delta\Gamma^{\mu}(E_2)} = 1 \quad \text{for all } E_1, E_2 \ll E_P
		\end{equation}
		
		\item \textbf{Cosmological scale ratios} \textbf{[S]}:
		\begin{equation}
			\frac{\kappa}{H_0} = \xi \approx \frac{\lambda_h^2 v^2}{16\pi^3 E_h^2}
		\end{equation}
		\textit{(Corrected on 1~Oct.~2026: previously $\xi = \lambda_h^2 v^2/(16\pi^3 E_h^2)$ -- the Higgs expression gives $\approx 1.30\times10^{-4}$, about $2.3\,\%$ below $4/30000$; a consistency check, not an equality.)}
	\end{enumerate}
	
	\section{Conclusions}
	
	The verification supports the central statement of the T0 model: \textbf{Fundamental physics can be based on scale ratios instead of assigned parameters}. The $\xi$ ratio characterizes these proportionalities and enables a description of physical phenomena with the single parameter $\xi$ (plus integer structural choices as motivated settings; SI anchors serve for conversion).
	
	\begin{thebibliography}{9}
		
		\bibitem{pascher_h0_energy_2025}
		Pascher, J. (2025). \textit{Pure Energy Formulation of $H_0$ and $\kappa$ Parameters in T0 Model Framework}. \\
		\url{https://github.com/jpascher/T0-Time-Mass-Duality/blob/main/2/pdf/xxx_H0_kappa_En.pdf}
		
		\bibitem{pascher_beta_derivation_2025}
		Pascher, J. (2025). \textit{Field Theoretical Derivation of $\beta_T$ Parameter in Natural Units ($\hbar = c = 1$)}. \\
		\url{https://github.com/jpascher/T0-Time-Mass-Duality/blob/main/2/pdf/093_DerivationVonBeta_En.pdf}
		
		\bibitem{pascher_elimination_mass_2025}
		Pascher, J. (2025). \textit{Elimination of Mass as Dimensional Placeholder in T0 Model: Toward Truly Parameter-Free Physics}. \\
		\url{https://github.com/jpascher/T0-Time-Mass-Duality/blob/main/2/pdf/052_EliminationOfMass_En.pdf}
		
		\bibitem{pascher_mol_candela_2025}
		Pascher, J. (2025). \textit{T0 Model: Universal Energy Relationships for Mole and Candela Units - Complete Derivation from Energy Scaling Principles}. \\
		\url{https://github.com/jpascher/T0-Time-Mass-Duality/blob/main/2/pdf/062_Moll_Candela_En.pdf}
		
	\end{thebibliography}
